\begin{abstract}
大型转动设备（汽轮机、发电机、压缩机等）的轴系对中精度直接决定机组的运行
平稳性与使用寿命。以火电厂汽轮机为例，转子直径通常在 1--2\,m 之间，对中精度
要求达到两丝（$20\,\mu\mathrm{m}$）量级。目前行业主流的钢丝找中法需要在待对中
轴之间反复拆装张紧的钢丝，不仅频繁干扰现场其他施工、在场地内构成绊拌等安全
隐患，钢丝自重引起的下坠挠度还会带来系统性的测量偏差。本文提出 DCPAM
（Dual-Camera-Based Point-to-Axis Measurement），用一条激光基准线替代钢丝，
通过前、后两个相机模块拍摄激光在毛玻璃接收屏上形成的光斑，经反投影、位姿对齐、
镜面虚像重建与点到直线距离计算，在统一的设备坐标系下恢复激光基准轴的三维空间
直线，并求解任意目标点到该轴的垂直距离。与只能读取激光斑在自身探测器上位置的
位置敏感探测器、需要合作反射器的激光跟踪仪，以及精度停留在毫米级的学术双目
激光线重建方法相比，DCPAM 直接面向“空间任意目标点到激光基准轴垂距”这一测量量，
目标精度为 $\pm15$--$50\,\mu\mathrm{m}$。实验部分给出重复精度评估与误差传播分析，
结果表明\todo{重复精度与主导误差来源结论}。

\end{abstract}

\begin{IEEEkeywords}
轴系对中，激光基准线，双目视觉，PnP 位姿估计，点到直线距离，镜面反射，精度分析。
\end{IEEEkeywords}
